% !TeX root = ../../../../main.tex
% chapters/sections/chapter5/subsections/a_shock/welfare.tex

\paragraph{厚生についての順位}
\label{sec:chapter5_a_shock_welfare}
次に価格分散\(\Delta_t^H\)を考慮した場合を例としてどの政策が最も厚生\(W_s^H\)の下落を抑制できているかを確認する。
厚生\(W_s^H\)の大きい順に政策を示す。
\begin{table}[H]
\centering
\caption{厚生についての順位}
\label{tab:chapter5_a_shock_welfare}
\begin{tabular}{lc}
\hline
金融政策 & \(W_s^H\)（\(\Delta_t^H\)有り） \\
\hline
実質消費水準目標（CLT） & $-1.3289$ \\
生産水準目標（OLT） & $-1.3289$ \\
名目総消費水準目標（NCLT） & $-1.3403$ \\
名目GDP水準目標（NGDPLT） & $-1.3403$ \\
消費者物価水準目標（CPLT） & $-1.3654$ \\
消費者物価インフレ目標（IT-CPI） & $-1.3657$ \\
生産者物価インフレ目標（IT-PPI） & $-1.4925$ \\
消費者物価テイラールール（TR-CPI） & $-1.5144$ \\
生産者物価テイラールール（TR-PPI） & $-1.5144$ \\
潜在生産水準目標（POLT） & $-1.5480$ \\
生産者物価水準目標（PPLT） & $-1.5486$ \\
\hline
\end{tabular}
\end{table}
CLTおよびOLTが最も高い厚生\(W_s^H\)を達成しており\(a\)ショックによる負の影響を最小限に留めていることがわかる。
一方で生産者物価指数（PPI）\(p_t^H\)や潜在生産\(y_t^{H,pot}\)のみを目標に含むPPLTやPOLTは極めて低い評価となった。